% =============================================================================
% GEMEINSAMES LAYOUT DER MECH-PILOT-KARTEN
% =============================================================================
% Diese Datei enthaelt nur das Design. Die Pilotendaten stehen getrennt in
% vorlage1.tex bis vorlage9.tex.

% Gemeinsamer Katalog aller Special Pilot Abilities
% =============================================================================
% BATTLETECH OVERRIDE - KATALOG DER SPECIAL PILOT ABILITIES
% =============================================================================
% Verwendung in mechpilot_card.tex:
%
%   \newcommand{\AbilityOne}{\RangeMasterLong}
%   \newcommand{\AbilityTwo}{\SpeedDemon}
%
% Der Kurzname steht jeweils im ersten Argument von \DeclareSPA und wird
% automatisch als LaTeX-Befehl angelegt. Leerzeichen und Sonderzeichen werden
% im Kurznamen weggelassen, zum Beispiel:
%
%   Range Master (Long) -> \RangeMasterLong
%   Eagle's Eyes        -> \EaglesEyes
%   Multi-Tasker        -> \MultiTasker
%
% Eine neue Fertigkeit wird nach diesem Muster ergänzt:
%
%   \DeclareSPA{Kurzname}{Slot}{Anzeigename}{Beschreibung}
% =============================================================================

\newcommand{\DeclareSPA}[4]{%
  \expandafter\def\csname #1\endcsname{#1}%
  \expandafter\def\csname SPAName#1\endcsname{#3}%
  \expandafter\def\csname SPAText#1\endcsname{#4}%
  \expandafter\def\csname SPASlot#1\endcsname{#2}%
}

\newcommand{\SPAName}[1]{%
  \ifcsname SPAName#1\endcsname
    \csname SPAName#1\endcsname
  \else
    Unbekannte SPA: #1%
  \fi
}

\newcommand{\SPAText}[1]{%
  \ifcsname SPAText#1\endcsname
    \csname SPAText#1\endcsname
  \else
    Für diese SPA wurde noch keine Beschreibung hinterlegt.%
  \fi
}

\newcommand{\SPASlot}[1]{%
  \ifcsname SPASlot#1\endcsname
    \csname SPASlot#1\endcsname
  \else
    ?%
  \fi
}

% ----------------------------------------------------------------------------
% AB SLOT 1
% ----------------------------------------------------------------------------

\DeclareSPA{BloodStalker}{1}{Blood Stalker}{%
At game start, designate one enemy. Attacks against it receive a -1 Target Number modifier; attacks against all other enemies suffer a +2 modifier. At the start of Movement, designate a new enemy if the current target is destroyed or out of line of sight.%
}

\DeclareSPA{ClusterHitter}{1}{Cluster Hitter}{%
As long as the unit with this SPA does not move during its Movement Phase, it can turn one M die to a 1 or one C die to a 3 for each TIC with M or C dice.%
}

\DeclareSPA{Dodge}{1}{Dodge}{%
Any opposing unit that attempts to deliver a physical attack against a unit with this SPA will suffer a +2 Target Number modifier.%
}

\DeclareSPA{EaglesEyes}{1}{Eagle's Eyes}{%
This unit adds 6 inches to all of its available sensor ranges.%
}

\DeclareSPA{EnvironmentalSpecialist}{1}{Environmental Specialist}{%
This unit reduces the mission's movement impairment by one level: Heavy$\rightarrow$Moderate, Moderate$\rightarrow$Light, Light$\rightarrow$None.%
}

\DeclareSPA{ForwardObserver}{1}{Forward Observer}{%
If the Forward Observer makes its own attack, any indirect attacks it spots for (IF or indirect Artillery) do not take the Target Number modifier for the spotter attacking.%
}

\DeclareSPA{ManeuveringAce}{1}{Maneuvering Ace}{%
This unit reduces the cost for moving through all woods and jungle terrain types by 1 inch per inch of movement.%
}

\DeclareSPA{MeleeMaster}{1}{Melee Master}{%
Add additional damage equal to half their unit's Size value (rounded up), upon delivering a successful physical attack of any kind, including punch, kick, melee, charging, and Death from Above attacks.%
}

\DeclareSPA{MeleeSpecialist}{1}{Melee Specialist}{%
A pilot with this ability applies an additional -1 Target Number modifier when making any physical attacks.%
}

\DeclareSPA{MultiTasker}{1}{Multi-Tasker}{%
This unit can ignore the target modifier for the first secondary target.%
}

\DeclareSPA{RangeMasterExtreme}{1}{Range Master (Extreme)}{%
The gunner for this unit specializes in attacks in extreme range: apply a -1 Target Number modifier for attacks at extreme range, but a +2 modifier for any attack made in short range.%
}

\DeclareSPA{RangeMasterLong}{1}{Range Master (Long)}{%
The gunner for this unit specializes in attacks in long range: apply a -1 Target Number modifier for attacks at long range, but a +2 modifier for any attack made in short range.%
}

\DeclareSPA{RangeMasterMedium}{1}{Range Master (Medium)}{%
The gunner for this unit specializes in attacks in medium range: apply a -1 Target Number modifier for attacks at medium range, but a +2 modifier for any attack made in short range.%
}

\DeclareSPA{StandAside}{1}{Stand-Aside}{%
This unit can move through hostile units during its Movement Phase, at an additional cost of 1 inch of Move. This action causes no damage to either unit.%
}

% ----------------------------------------------------------------------------
% AB SLOT 2
% ----------------------------------------------------------------------------

\DeclareSPA{Hopper}{2}{Hopper}{%
A unit controlled by a pilot with this SPA can ignore the effects of the first MP Hit it receives.%
}

\DeclareSPA{JumpingJack}{2}{Jumping Jack}{%
A pilot with the Jumping Jack SPA is so comfortable with the use of jumping movement that their unit uses a +1 attacker movement modifier for jumping instead of +2.%
}

\DeclareSPA{Marksman}{2}{Marksman}{%
As long as this unit stands still in the Movement Phase, it can choose the following for its first successful weapon hit: reduce the damage of the weapon by half (rounded down, to a minimum of 1). If it does, it performs an additional critical hit check, even if the location is still armored.%
}

\DeclareSPA{SpeedDemon}{2}{Speed Demon}{%
A pilot with the Speed Demon SPA adds +2 inches to the unit's ground movement and +3 inches to the unit's sprint value, without increasing the TMM.%
}

\DeclareSPA{Swordsman}{2}{Swordsman}{%
On a successful physical attack made while using the unit's melee weapon (axe, sword etc.), this pilot may choose one of two options: deliver 2 extra points of damage to the opponent, or roll one additional Critical Hit against the target, even if it still has armor remaining. Note that if the unit lacks a melee weapon, the Swordsman SPA will have no effect.%
}

% ----------------------------------------------------------------------------
% AB SLOT 3
% ----------------------------------------------------------------------------

\DeclareSPA{HotDog}{3}{Hot Dog}{%
This unit acts as if it was one level lower on the heat scale than it actually is, except when hitting automatic shutdown.%
}

\DeclareSPA{NaturalGrace}{3}{Natural Grace}{%
This unit may make attacks as if it has a 360-degree firing arc, but is still hit on the rear location table if an attack originates from the rear arc.%
}

\DeclareSPA{ObliqueAttacker}{3}{Oblique Attacker}{%
This unit receives a -1 Target Number modifier for indirect attacks and may even make indirect fire attacks without a friendly spotter. If attempting to use indirect fire without a friendly spotter, however, the unit trades its -1 Target Number modifier for a +2 modifier, which replaces any and all spotter-related modifiers.%
}

\DeclareSPA{Sniper}{3}{Sniper}{%
This gunner's SPA reduces their Range Modifiers at medium, long and extreme range by -1.%
}


% Farben, angelehnt an die Override-Oberflaeche
\definecolor{CardBlack}{HTML}{0C0D0D}
\definecolor{Panel}{HTML}{171919}
\definecolor{PanelLight}{HTML}{202323}
\definecolor{LineGray}{HTML}{343838}
\definecolor{TextMain}{HTML}{F2F2F2}
\definecolor{TextMuted}{HTML}{A7AAAA}
\definecolor{OverrideGreen}{HTML}{6DDF73}
\definecolor{OverrideDarkGreen}{HTML}{183D20}

\newcommand{\SmallCapsLabel}[1]{%
  {\fontsize{4.7}{5.2}\selectfont\bfseries\color{TextMuted}\MakeUppercase{#1}}%
}

\newcommand{\StatBox}[4]{%
  % #1 x links, #2 Breite, #3 Label, #4 Wert
  \fill[Panel] (#1,56.2) rectangle ++(#2,11.4);
  \fill[OverrideGreen] (#1,56.2) rectangle ++(0.55,11.4);
  \node[anchor=north west, text=TextMuted, font=\fontsize{4.2}{4.8}\selectfont\bfseries]
    at ($(#1,67.6)+(1.7,-1.8)$) {\MakeUppercase{#3}};
  \node[anchor=west, text=TextMain, font=\fontsize{13}{13}\selectfont\bfseries]
    at ($(#1,56.2)+(2.1,4.2)$) {#4};
}

\newcommand{\AbilityBox}[5]{%
  % #1 unteres y, #2 oberes y, #3 Nummer, #4 Name, #5 Text
  \fill[Panel] (3,#1) rectangle (60.5,#2);
  \draw[LineGray, line width=0.25mm] (3,#1) rectangle (60.5,#2);
  \fill[OverrideGreen] (3,#1) rectangle (3.65,#2);
  \fill[OverrideDarkGreen] (5.1,{#2-4.7}) rectangle (9.4,{#2-1.0});
  \node[text=OverrideGreen, font=\fontsize{5.5}{6}\selectfont\bfseries]
    at (7.25,{#2-2.85}) {#3};
  \node[anchor=north west, text=TextMain,
        font=\fontsize{5.3}{5.8}\selectfont\bfseries]
    at (10.6,{#2-1.0}) {\MakeUppercase{#4}};
  % Lange Beschreibungen werden automatisch verkleinert.
  \node[anchor=north west, inner sep=0pt, text=TextMuted]
    at (10.6,{#2-4.8}) {%
      \begin{adjustbox}{max totalsize={47.7mm}{8.2mm}}
        \begin{minipage}{47.7mm}
          \RaggedRight\fontsize{4.35}{5.05}\selectfont #5
        \end{minipage}
      \end{adjustbox}%
    };
}

% Erzeugt genau eine Karte im Format 63,5 x 88,9 mm.
% Die benoetigten Variablen werden in der jeweiligen vorlageX.tex gesetzt.
\newcommand{\MechPilotCard}{%
\begin{tikzpicture}[x=1mm,y=1mm]
  % Exakte Bounding Box, damit die Anordnung auf A4 masshaltig bleibt.
  \path[use as bounding box] (0,0) rectangle (63.5,88.9);

  % Kartenhintergrund und Aussenrand
  \clip[rounded corners=2.4mm] (0.4,0.4) rectangle (63.1,88.5);
  \fill[CardBlack] (0,0) rectangle (63.5,88.9);
  \draw[OverrideGreen, line width=0.55mm, rounded corners=2.2mm]
    (1.0,1.0) rectangle (62.5,87.9);

  % Kopfbereich
  \fill[Panel] (2.4,69.6) rectangle (61.1,86.8);
  \fill[OverrideGreen] (2.4,69.6) rectangle (3.1,86.8);

  % Profilbild mit Platzhalter, wenn die Datei fehlt
  \fill[CardBlack] (5.0,72.0) rectangle (19.4,84.9);
  \begin{scope}
    \clip[rounded corners=0.7mm] (5.0,72.0) rectangle (19.4,84.9);
    \IfFileExists{\PilotImage}{%
      \node[inner sep=0pt] at (12.2,78.45)
        {\includegraphics[width=14.4mm,height=12.9mm]{\PilotImage}};
    }{%
      \fill[PanelLight] (5.0,72.0) rectangle (19.4,84.9);
      \node[text=TextMuted, align=center,
            font=\fontsize{4.2}{4.8}\selectfont\bfseries]
        at (12.2,78.45) {PROFILBILD\\EINFUEGEN};
    }
  \end{scope}
  \draw[OverrideGreen, line width=0.35mm, rounded corners=0.7mm]
    (5.0,72.0) rectangle (19.4,84.9);

  % Status und Name
  \fill[OverrideDarkGreen] (22.1,82.0) rectangle (31.0,84.3);
  \node[text=OverrideGreen, font=\fontsize{4.2}{4.5}\selectfont\bfseries]
    at (26.55,83.15) {\PilotStatus};

  \fill[PanelLight] (31.8,82.0) rectangle (43.2,84.3);
  \node[text=TextMuted, font=\fontsize{4.0}{4.4}\selectfont\bfseries]
    at (37.5,83.15) {\PilotRole};

  \node[anchor=west, text=TextMain,
        font=\fontsize{12.5}{13}\selectfont\bfseries]
    at (22.0,77.4) {``\PilotName''};

  \node[anchor=west, text=TextMuted,
        font=\fontsize{4.6}{5}\selectfont\bfseries]
    at (22.1,73.4) {BATTLETECH OVERRIDE};

  % Kampfwerte
  \node[anchor=west] at (3.0,68.6) {\SmallCapsLabel{Kampfwerte}};
  \draw[LineGray, line width=0.25mm] (3.0,68.0) -- (60.5,68.0);

  \StatBox{3.0}{18.3}{Gunnery}{\Gunnery}
  \StatBox{22.6}{18.3}{Piloting}{\Piloting}
  \StatBox{42.2}{18.3}{SPA-Slots}{\SPASlots}

  % Spezialfaehigkeiten
  \node[anchor=west] at (3.0,54.8) {\SmallCapsLabel{Special Pilot Abilities}};
  \draw[LineGray, line width=0.25mm] (3.0,54.2) -- (60.5,54.2);

  \AbilityBox{39.5}{52.8}{1}{\SPAName{\AbilityOne}}{\SPAText{\AbilityOne}}
  \AbilityBox{24.4}{37.7}{2}{\SPAName{\AbilityTwo}}{\SPAText{\AbilityTwo}}

  % Einheitenzuweisung
  \node[anchor=west] at (3.0,22.7) {\SmallCapsLabel{Einheitenzuweisung}};
  \draw[LineGray, line width=0.25mm] (3.0,22.1) -- (60.5,22.1);

  \fill[Panel] (3.0,7.4) rectangle (60.5,20.7);
  \draw[LineGray, line width=0.25mm] (3.0,7.4) rectangle (60.5,20.7);
  \fill[OverrideDarkGreen] (5.2,12.2) rectangle (10.4,17.4);
  \node[text=OverrideGreen, font=\fontsize{5.5}{6}\selectfont\bfseries]
    at (7.8,14.8) {BM};
  \node[anchor=west, text=TextMain,
        font=\fontsize{7.2}{7.8}\selectfont\bfseries]
    at (12.2,16.1) {\AssignedUnit};
  \node[anchor=west, text=TextMuted,
        font=\fontsize{5.0}{5.5}\selectfont]
    at (12.2,12.4) {\AssignedVariant};

  % Fusszeile
  \node[anchor=south east, text=TextMuted,
        font=\fontsize{3.8}{4.2}\selectfont\bfseries]
    at (60.5,3.0) {MECH PILOT CARD};

\end{tikzpicture}%
}
